\section{实验目的}

PA4 的目标是在 PA3 已经完成的异常控制流、系统调用、文件系统和设备文件基础上，进一步理解操作系统如何支持多个用户程序共同运行。本阶段的关键词是“虚拟地址空间”和“分时多任务”：前者让不同用户进程可以使用相同的虚拟入口和相似的运行布局而互不覆盖，后者让操作系统通过上下文切换把处理器时间分配给多个进程。

从系统层次看，本实验跨越 NEMU、AM、Nanos-lite 和 Navy-apps 四层。NEMU 提供分页、异常和硬件中断的模拟；AM 把体系结构相关的 trap frame、页表和中断入口包装成统一接口；Nanos-lite 负责进程地址空间、PCB、调度和系统调用处理；Navy-apps 则作为用户程序运行在上述机制之上。PA4 的核心训练不只是补全若干函数，而是理解这些层次怎样共同支撑仙剑奇侠传、hello 和 videotest 的分时运行。

\section{实验内容与阶段划分}

按照实验指导书，PA4 可以划分为三个主要阶段和一个展示任务。第一阶段实现分页和用户虚拟地址空间；第二阶段实现 trap frame 构造、上下文切换和基于系统调用返回的协作式分时；第三阶段加入时钟中断，让进程切换不再依赖用户程序主动发起系统调用；最后通过 F12 在 PAL 与 videotest 之间切换，展示完整系统。

\begin{table}[H]
\centering
\caption{PA4 实验内容与阶段划分}
\label{tab:pa4-stage-plan}
\small
\begin{tabularx}{\textwidth}{p{0.16\textwidth} p{0.23\textwidth} X p{0.22\textwidth}}
\toprule
阶段 & 目标 & 主要内容 & 验证方式 \\
\midrule
阶段一 & 建立用户地址空间 & NEMU 支持分页地址转换；AM 初始化内核页表并实现 \texttt{\_protect()}、\texttt{\_map()}、\texttt{\_switch()}；loader 将用户程序加载到虚拟地址空间 & 开启 \texttt{HAS\_PTE} 后用户程序仍能运行 \\
\midrule
阶段二 & 实现上下文切换 & 实现 \texttt{\_trap()}、\texttt{\_umake()} 和 \texttt{schedule()}；修改 \texttt{asm\_trap} 根据返回的 trap frame 切换栈；加载 PAL 与 hello & PAL 和 hello 能交替运行 \\
\midrule
阶段三 & 实现抢占式分时 & NEMU 添加时钟中断；AM 将中断包装为 \texttt{\_EVENT\_IRQ\_TIME}；Nanos-lite 在时钟中断中调度 & 调度不再依赖系统调用，死循环进程不能长期独占 CPU \\
\midrule
展示任务 & 切换图形程序 & 加载 \texttt{/bin/videotest}；用 \texttt{current\_game} 记录当前图形进程；在 \texttt{events\_read()} 中响应 F12 & F12 在 PAL 和 videotest 间切换 \\
\bottomrule
\end{tabularx}
\end{table}

\section{阶段一：分页与用户虚拟地址空间}

\subsection{分页机制的基本路径}

开启分页后，CPU 发出的虚拟地址不再直接作为物理地址使用，而要经过页目录和页表两级转换。NEMU 中需要通过 \texttt{CR0.PG} 判断分页是否开启，通过 \texttt{CR3} 找到当前页目录。如果访问没有跨页，只需把虚拟地址拆成页目录索引、页表索引和页内偏移，再查出物理页框；如果访问跨越页边界，则必须拆成两次访问分别转换。这个细节会影响字符串、指令取数和用户缓冲区读写，是分页阶段容易隐藏错误的位置。

AM 的 \texttt{\_pte\_init()} 首先建立内核恒等映射，即虚拟地址和物理地址在内核区保持一致。这样打开分页后，Nanos-lite 的代码、数据、设备映射和内核栈仍然可以继续使用原来的地址访问。随后每个用户进程通过 \texttt{\_protect()} 得到自己的页目录，该页目录复制内核映射，再在用户区加入本进程独有的用户页映射。

\subsection{控制寄存器初始化}

分页机制依赖 x86 的控制寄存器。其中 \texttt{CR0} 记录处理器是否处于保护模式、是否开启分页等状态，\texttt{CR3} 记录当前页目录的物理基地址。NEMU 在 \texttt{CPU\_state} 中保存这两个寄存器的软件表示；在复位函数 \texttt{restart()} 中，需要给它们设置初始值。实验中将 \texttt{CR0} 初始化为 \texttt{0x60000011}，使其具备 PA4 指导书要求的初始控制状态；\texttt{CR3} 初始化为 0，避免模拟器重启后残留旧页目录基地址。

\begin{lstlisting}[language=C,caption={\texttt{CPU\_state} 中的 CR0/CR3 寄存器字段},frame=single]
typedef struct {
  union {
    union {
      uint32_t _32;
      uint16_t _16;
      uint8_t _8[2];
    } gpr[8];

    struct {
      rtlreg_t eax, ecx, edx, ebx, esp, ebp, esi, edi;
    };
  };

  vaddr_t eip;
  rtlreg_t cs;

  union {
    struct {
      uint32_t CF       : 1;
      uint32_t          : 5;
      uint32_t ZF       : 1;
      uint32_t SF       : 1;
      uint32_t          : 1;
      uint32_t IF       : 1;
      uint32_t          : 1;
      uint32_t OF       : 1;
      uint32_t          : 19;
    };
    rtlreg_t eflags;
  };

  CR0 cr0;
  CR3 cr3;

  struct {
    uint16_t limit;
    paddr_t base;
  } idtr;
} CPU_state;
\end{lstlisting}

\begin{lstlisting}[language=C,caption={\texttt{restart()} 中初始化 CR0/CR3},frame=single]
static inline void restart() {
  /* Set the initial instruction pointer. */
  cpu.eip = ENTRY_START;
  cpu.cs = 8;
  cpu.eflags = 0x00000002;
  cpu.cr0.val = 0x60000011;
  cpu.cr3.val = 0;

#ifdef DIFF_TEST
  init_qemu_reg();
#endif
}
\end{lstlisting}

上述代码使 NEMU 的初始处理器状态与后续 PA4 分页实验保持一致。之后 AM 在 \texttt{\_pte\_init()} 中调用 \texttt{set\_cr3(kpdirs)} 设置内核页目录，并通过 \texttt{set\_cr0(get\_cr0() | CR0\_PG)} 打开分页位。这样，真正进入分页模式的时机仍由 AM 控制，而 NEMU 的复位状态已经具备正确的控制寄存器基础。

\subsection{虚拟地址读写与页表翻译}

在 PA3 中，\texttt{vaddr\_read()} 和 \texttt{vaddr\_write()} 直接把虚拟地址当作物理地址使用。PA4 开启分页后，这两个接口必须先检查 \texttt{CR0.paging}：如果分页未开启，继续沿用直接访问；如果分页已经开启，就通过当前 \texttt{CR3} 指向的页目录完成二级页表翻译。当前实现暂时不处理跨页访问，一旦一次读写跨越页边界就触发 \texttt{assert(0)}，后续可以再扩展为拆分成两次物理访问。

\subsubsection*{\texttt{page\_translate()} 的实现逻辑}

\texttt{page\_translate()} 的输入是一个 32 位虚拟地址，输出是翻译后的物理地址。函数首先按照 4KB 页大小拆分虚拟地址：高 10 位是页目录索引 \texttt{dir}，中间 10 位是页表索引 \texttt{page}，低 12 位是页内偏移 \texttt{offset}。然后，函数通过 \texttt{cpu.cr3.page\_directory\_base << 12} 得到当前页目录的物理基地址，并从页目录中读取对应的 PDE。若 PDE 的 \texttt{present} 位没有置位，说明该虚拟地址所在的大区间没有有效页表，直接触发断言。

在 PDE 有效的情况下，函数继续根据 \texttt{pde.page\_frame << 12} 找到页表物理基地址，再读取对应的 PTE。PTE 的 \texttt{present} 位表示最终虚拟页是否已经映射到物理页框；只有该位有效时，翻译才可以继续。最后函数把 PTE 中的物理页框号左移 12 位，并与原虚拟地址的页内偏移合并，得到完整物理地址。整个过程中读取页目录项和页表项都使用 \texttt{paddr\_read()}，因为页表本身存放在物理内存中，不应该再次走虚拟地址翻译。

\subsubsection*{\texttt{vaddr\_read()} 的实现逻辑}

\texttt{vaddr\_read()} 是所有虚拟地址读访问的统一入口。函数首先检查 \texttt{cpu.cr0.paging}：如果分页位没有打开，说明当前仍处于 PA3 式的直通内存模型，虚拟地址可以直接交给 \texttt{paddr\_read()} 读取。这样做可以保证开启分页前的 NEMU 启动、加载镜像和早期代码执行路径不受影响。

如果分页已经开启，函数接着判断本次读取是否跨越页边界。判断条件是 \texttt{(addr \& PAGE\_MASK) + len > PAGE\_SIZE}，其中 \texttt{addr \& PAGE\_MASK} 表示页内偏移。如果页内偏移加读取长度超过 4KB，说明一次访问覆盖了两个虚拟页；当前阶段先用 \texttt{assert(0)} 暴露这种情况。非跨页时，函数调用 \texttt{page\_translate()} 得到物理地址，再复用 \texttt{paddr\_read()} 完成实际读取。

\subsubsection*{\texttt{vaddr\_write()} 的实现逻辑}

\texttt{vaddr\_write()} 与 \texttt{vaddr\_read()} 的结构对称，但它负责写入。分页未开启时，它直接调用 \texttt{paddr\_write(addr, len, data)}，保持原有直通行为，并立即返回。分页开启后，它使用同样的跨页判断条件检查本次写入是否落在单个页内；如果跨页，当前先断言失败，避免在未实现拆分写入的情况下错误覆盖内存。

当写入范围没有跨页时，函数调用 \texttt{page\_translate()} 将虚拟起始地址翻译为物理地址，再调用 \texttt{paddr\_write()} 写入数据。由于最终仍然走物理写接口，因此 MMIO 写入、普通物理内存写入等底层行为仍由 \texttt{paddr\_write()} 统一处理，虚拟地址层只负责“是否分页”和“如何翻译地址”。

\begin{lstlisting}[language=C,caption={\texttt{memory.c} 中虚拟地址翻译和读写实现},frame=single]
static paddr_t page_translate(vaddr_t addr) {
  uint32_t dir = (addr >> 22) & 0x3ff;
  uint32_t page = (addr >> 12) & 0x3ff;
  uint32_t offset = addr & PAGE_MASK;

  paddr_t pdir_base = cpu.cr3.page_directory_base << 12;
  PDE pde;
  pde.val = paddr_read(pdir_base + dir * sizeof(PDE), 4);
  assert(pde.present);

  paddr_t ptab_base = pde.page_frame << 12;
  PTE pte;
  pte.val = paddr_read(ptab_base + page * sizeof(PTE), 4);
  assert(pte.present);

  return (pte.page_frame << 12) | offset;
}

uint32_t vaddr_read(vaddr_t addr, int len) {
  if (!cpu.cr0.paging) {
    return paddr_read(addr, len);
  }

  if ((addr & PAGE_MASK) + len > PAGE_SIZE) {
    assert(0);
  }

  paddr_t paddr = page_translate(addr);
  return paddr_read(paddr, len);
}

void vaddr_write(vaddr_t addr, int len, uint32_t data) {
  if (!cpu.cr0.paging) {
    paddr_write(addr, len, data);
    return;
  }

  if ((addr & PAGE_MASK) + len > PAGE_SIZE) {
    assert(0);
  }

  paddr_t paddr = page_translate(addr);
  paddr_write(paddr, len, data);
}
\end{lstlisting}

三者的分工可以概括为：\texttt{page\_translate()} 只负责页表查询和地址转换；\texttt{vaddr\_read()} 负责读路径上的分页开关判断、跨页检查和物理读调用；\texttt{vaddr\_write()} 负责写路径上的分页开关判断、跨页检查和物理写调用。这样的划分让页表翻译逻辑集中在一个函数中，同时保留了原有物理内存接口对 MMIO 和普通内存的统一处理。

\subsection{用户进程地址空间}

每个 PCB 中保存一个 \texttt{\_Protect as}，它描述该进程的虚拟地址空间。创建进程时，\texttt{\_protect()} 为进程分配页目录并复制内核页目录；加载程序时，loader 按页分配物理页，用 \texttt{\_map()} 将用户虚拟地址映射到这些物理页，然后把可执行文件内容读入对应物理页。这样，PAL 和 hello 可以都从同一个虚拟入口运行，但它们的虚拟页最终指向不同的物理页，不会互相覆盖。

\begin{figure}[H]
\centering
\begin{tikzpicture}[node distance=1.2cm, >=Latex, font=\small]
  \node[draw, rounded corners, minimum width=3.2cm, minimum height=0.8cm] (palva) {PAL 虚拟地址 0x4000000};
  \node[draw, rounded corners, right=2.4cm of palva, minimum width=3.2cm, minimum height=0.8cm] (palpa) {物理页 A};
  \node[draw, rounded corners, below=1.0cm of palva, minimum width=3.2cm, minimum height=0.8cm] (hellova) {hello 虚拟地址 0x4000000};
  \node[draw, rounded corners, right=2.4cm of hellova, minimum width=3.2cm, minimum height=0.8cm] (hellopa) {物理页 B};
  \draw[->] (palva) -- node[above]{PAL 页表} (palpa);
  \draw[->] (hellova) -- node[below]{hello 页表} (hellopa);
\end{tikzpicture}
\caption{不同进程可使用相同虚拟地址而映射到不同物理页}
\label{fig:pa4-address-space}
\end{figure}

堆区管理也需要纳入地址空间机制。PA3 中 \texttt{SYS\_brk} 可以简单返回成功；PA4 中，\texttt{mm\_brk()} 需要根据当前进程的 \texttt{cur\_brk} 和 \texttt{max\_brk} 判断是否要分配新物理页并建立映射。这样用户态 \texttt{malloc()} 扩展堆区时，内核会把相应虚拟页纳入当前进程的地址空间。

\section{阶段二：上下文切换与协作式分时}

\subsection{从函数调用到 trap frame 恢复}

PA3 中 Nanos-lite 可以在 loader 返回后直接把入口地址转换成函数指针并调用用户程序。但在多任务系统中，这种方式不合适：用户程序的现场应当由异常返回机制恢复，而不是由内核直接调用。PA4 中需要用 \texttt{\_umake()} 在用户栈上人工构造 trap frame，让第一次调度时可以像从一次异常处理返回一样进入用户程序。

\begin{lstlisting}[language=C,caption={用户进程初始现场的关键字段},frame=single]
tf->eip = (uintptr_t)entry;
tf->cs = 8;
tf->eflags = 0x2;       /* later set IF for timer interrupt */
tf->eax = tf->ebx = tf->ecx = tf->edx = 0;
\end{lstlisting}

由于 Navy-apps 的入口是 \texttt{\_start()}，而 \texttt{\_start()} 会按照普通 C 运行时约定读取参数，因此还需要在 trap frame 上方准备一个简单的初始栈帧，将 \texttt{argc}、\texttt{argv} 和 \texttt{envp} 设置为 0 或 \texttt{NULL}。用户程序不会从 \texttt{\_start()} 返回，因此返回地址不需要具备实际意义。

\subsection{调度器与栈切换}

上下文切换的本质是切换 trap frame 所在的栈。发生异常或中断时，当前进程的寄存器、EIP、CS、EFLAGS 等信息被保存到当前栈上；Nanos-lite 的 \texttt{schedule()} 保存旧进程的 trap frame 指针，选择新进程，切换到新进程的页目录，然后返回新进程的 trap frame 指针。AM 的 \texttt{asm\_trap} 收到该返回值后，将 \texttt{\%esp} 改为新 trap frame 的地址，之后的 \texttt{popal} 和 \texttt{iret} 就会恢复新进程的现场。

\begin{lstlisting}[language={[x86masm]Assembler},caption={asm\_trap 中完成上下文切换的关键动作},frame=single]
pushl %esp
call irq_handle
addl $4, %esp
movl %eax, %esp
popal
addl $8, %esp
iret
\end{lstlisting}

最初可以通过内核自陷 \texttt{\_trap()} 触发第一次调度，让 Nanos-lite 从初始化过程切换到第一个用户进程。随后为了验证多进程机制，可以在系统调用处理结束后调用 \texttt{schedule()}，让 PAL 与 hello 在每次系统调用返回前交替运行。这种方式属于协作式分时，因为切换依赖用户程序主动陷入内核。

\section{阶段三：硬件时钟中断与抢占式分时}

\subsection{为什么需要硬件中断}

如果调度只发生在系统调用返回前，那么一个陷入死循环且不发起系统调用的程序就可以永久占用 CPU。PA4 通过时钟中断解决这个问题：时钟设备定期向 CPU 发出中断请求，CPU 在允许中断的情况下暂停当前指令流，保存现场并进入中断处理程序。这样调度时机不再由用户程序决定，而由外部硬件事件决定。

在 NEMU 中，时钟设备的 \texttt{timer\_intr()} 周期性调用 \texttt{dev\_raise\_intr()}。更符合硬件语义的做法是让 \texttt{dev\_raise\_intr()} 只拉高 \texttt{cpu.INTR}，CPU 在每条指令执行结束后检查 \texttt{INTR} 和 \texttt{EFLAGS.IF}。如果二者都有效，就清除 \texttt{INTR} 并调用 \texttt{raise\_intr(0x20, cpu.eip)}，通过 IDT 跳转到 AM 的时钟中断入口。

\begin{lstlisting}[language=C,caption={时钟中断响应的逻辑},frame=single]
update_eip();

if (cpu.INTR && cpu.IF) {
  cpu.INTR = false;
  raise_intr(0x20, cpu.eip);
}
\end{lstlisting}

\subsection{从时钟中断到进程调度}

AM 在 IDT 中为 \texttt{0x20} 设置入口，并在 \texttt{irq\_handle()} 中把该中断号包装成 \texttt{\_EVENT\_IRQ\_TIME}。Nanos-lite 收到该事件后调用 \texttt{schedule()}，选择下一个进程。由于 \texttt{\_umake()} 构造的初始 \texttt{EFLAGS} 中设置了 IF 位，用户程序运行时能够响应外部时钟中断；而 \texttt{raise\_intr()} 在保存旧 \texttt{EFLAGS} 后清除 IF 位，可以避免普通中断处理过程中继续嵌套响应同级中断。

加入时钟中断后，原先“系统调用后调度”的临时代码可以移除。此时 PAL 和 hello 的切换由时间片驱动：无论当前运行的是输出字符串的 hello，还是计算和绘图较多的 PAL，只要时钟中断到达，Nanos-lite 就有机会重新获得控制权并切换进程。

\section{展示任务：F12 切换图形程序}

最终展示任务是在 PAL 和 videotest 之间切换。因为两个图形程序如果同时参与调度，会互相覆盖帧缓冲内容，所以只允许一个图形程序与 hello 分时运行。实现上可以在 Nanos-lite 中维护一个 \texttt{current\_game} 变量，初始指向 PAL；调度器在 \texttt{current\_game} 和 hello 之间选择进程。设备文件 \texttt{/dev/events} 的 \texttt{events\_read()} 读取到 \texttt{kd F12} 时切换 \texttt{current\_game}，于是下一轮调度就会切换到另一个图形程序。

\section{必答题：分页机制和硬件中断如何支撑分时运行}

分页机制支撑分时运行的关键在于隔离和重定位。PAL 和 hello 都可以使用相同的用户虚拟入口地址和类似的运行时布局，但它们的页目录不同。loader 为每个进程分配不同的物理页，并通过 \texttt{\_map()} 建立用户虚拟页到物理页的映射；调度器切换进程时调用 \texttt{\_switch()} 写入新的 \texttt{CR3}，于是同一个虚拟地址在不同进程中会被解释成不同的物理地址。与此同时，各进程页目录都保留内核映射，因此发生系统调用或中断后，CPU 虽然处在某个用户地址空间中，仍然能够访问 Nanos-lite 和 AM 的代码与数据。

硬件中断支撑分时运行的关键在于夺回控制权。NEMU 的时钟设备周期性设置中断请求，CPU 在 \texttt{IF=1} 时响应请求，通过 IDT 进入 AM 的中断入口。AM 保存当前进程现场并构造 trap frame，将中断封装为 \texttt{\_EVENT\_IRQ\_TIME} 交给 Nanos-lite。Nanos-lite 的 \texttt{schedule()} 保存当前进程的 trap frame，选择下一个进程并切换地址空间；随后 \texttt{asm\_trap} 把栈顶切换到新进程的 trap frame，最后由 \texttt{iret} 恢复新进程的 EIP、CS 和 EFLAGS。于是从用户程序视角看，自己只是被暂停和恢复；从操作系统视角看，处理器时间被切分给了多个进程。

因此，分页解决“多个进程的内存如何互不干扰”，硬件中断解决“操作系统如何定期重新获得 CPU”。二者结合后，PAL 和 hello 才能在同一套 NEMU、AM、Nanos-lite 系统上以分时方式运行。

\section{思考题：固定地址保存中断现场的灾难性后果}

如果硬件总是把中断信息固定保存在内存地址 \texttt{0x1000}，AM 也总是从这里构造 trap frame，那么一旦发生中断嵌套，第二次中断会覆盖第一次中断保存的现场。第一次中断原本需要依靠这些现场恢复被中断程序的 EIP、CS、EFLAGS 和通用寄存器；覆盖发生后，返回路径会读取到第二次中断的现场或者混杂后的错误现场。

这种错误会表现为控制流和数据状态同时损坏。轻则系统调用返回值、通用寄存器或栈指针异常，导致用户程序行为错乱；重则 \texttt{iret} 恢复到错误 EIP，CPU 跳转到无效地址或错误函数中继续执行，引发非法指令、页故障、死循环或内核崩溃。更严重的是，若两个进程的现场互相覆盖，还会出现进程串扰：一个进程可能带着另一个进程的寄存器和返回地址继续运行。使用栈保存现场的意义正在于此：每次异常和中断都会自然形成新的 trap frame，嵌套时也不会覆盖旧现场。

\section{实验总结}

PA4 将 PA3 中已经打通的异常控制流推广为完整的多进程运行机制。分页让每个用户进程拥有独立虚拟地址空间，上下文切换让处理器能够从一个 trap frame 恢复到另一个 trap frame，时钟中断则让调度不依赖用户程序的自觉行为。经过这些机制的组合，Nanos-lite 从“能运行一个用户程序的内核”推进到了“能让多个用户程序分时运行的操作系统雏形”。

本阶段最重要的认识是：分时多任务并不是某一个函数单独完成的功能，而是硬件、抽象机器和操作系统之间的协作结果。NEMU 模拟了分页和中断这样的硬件事实，AM 把事实包装成统一接口，Nanos-lite 决定调度策略和地址空间切换，用户程序则在被隔离出的虚拟世界中继续执行。正是这些层次的边界清晰又彼此衔接，才让 PAL、hello 和 videotest 能够共享同一个处理器而不互相破坏。
